\section{Energy-momentum tensor}\label{sec:emt}

In a continuous $D$-dimensional space-time, the gauge invariant part of the
\emt\ reads
\begin{equation}
  \begin{split}
    T_{\mu\nu}(x) \equiv \frac{1}{g_0^2}\left[
      \calo_{1,\mu\nu}(x)-\frac{1}{4}\calo_{2,\mu\nu}(x)\right]
    + \frac{1}{4} \calo_{3,\mu\nu}(x)
    - \frac{1}{2} \calo_{4,\mu\nu}(x)
    - \calo_{5 ,\mu\nu}(x) \,,
    \label{eq:emt}
  \end{split}
\end{equation}
where $g_0$ is the bare coupling constant of \qcd. The operators are
defined as
\begin{equation}
  \begin{split}
    \calo_{1,\mu\nu}(x) &\equiv
    F_{\mu\rho}^a(x)F_{\nu\rho}^a(x)\,,\\[10pt] \calo_{2,\mu\nu}(x)
    &\equiv
    \delta_{\mu\nu}F_{\rho\sigma}^a(x)F_{\rho\sigma}^a(x)\,,\\[10pt]
    \calo_{3_f,\mu\nu}(x) &\equiv
    \bar\psi_{f}(x)\left(\gamma_\mu\overleftrightarrow{D}\!_\nu
    +\gamma_\nu\overleftrightarrow{D}\!_\mu\right)\psi_{f}(x)\,,\\
    \calo_{4_f,\mu\nu}(x)
    &\equiv \delta_{\mu\nu}\bar\psi_{f}(x)
    \overleftrightarrow{\slashed{D}}\psi_{f}(x)\,,\phantom{\bigg(\bigg)}\\
    \calo_{5_f,\mu\nu}(x)
    &\equiv \delta_{\mu\nu}
    m_{f,0}\bar\psi_{f}(x)\psi_{f}(x)\,,\\
    \calo_{i,\mu\nu}(x) &=
    \sum_{f=1}^{n_F}\calo_{i_f,\mu\nu}(x)\,,\qquad i\in\{3,4,5\}\,,
    \label{eq:ops}
  \end{split}
\end{equation}
where $f$ labels the $n_F$ different quark flavors, $m_{f,0}$ is the
bare quark mass, and
\begin{equation}
  \begin{split}
    F^a_{\mu\nu} = \partial_\mu A_\nu^a - \partial_\nu A_\mu^a +
    f^{abc}A_\mu^bA_\nu^c\,,\qquad
    \overleftrightarrow{D}\!_\mu = \partial_\mu -
    \overleftarrow{\partial}\!_\mu + 2A_\mu\,.
  \end{split}
\end{equation}
The notation $i_f\in\{i_1,\ldots,i_{n_F}\}$ for the indices which label
different flavors will be useful later on in this paper.  In general,
$T_{\mu\nu}$ may contain gauge-dependent operators\, which vanish when
evaluating physical matrix elements \cite{Nielsen:1977sy}. Here and in
what follows, we implicitly assume that the vacuum expectations values
of all composite operators have been subtracted\footnote{In other words,
  the precise definition of $\calo_{1,\mu\nu}$, for example, would be
  given by $F_{\mu\rho}^aF_{\rho\nu}^a-\langle
  F_{\mu\rho}^aF_{\rho\nu}^a\rangle$.} so that $\langle
\calo_{i,\mu\nu}(x)\rangle\equiv 0\ \forall i$.

In this paper, we will focus on the case where physical matrix elements
of the \emt\ itself are considered, i.e.\ no other operator multiplies
the \emt\ at the same space-time point. In this case, the
equations-of-motion (\eom{}) render the set of operators in \eqn{eq:ops}
redundant. In particular, the \eom\ for the quark fields in regular
\qcd\ implies
\begin{equation}
  0 = \calo_{4, \mu \nu} (x) + 2 \, \calo_{5, \mu \nu} (x) \,,
  \label{eq:eom}
\end{equation}
which allows us to eliminate $\calo_{5,\mu\nu}$ from the set of
operators in \eqn{eq:ops}. Note that due to this relation, the last two
terms in \eqn{eq:emt} cancel.

We will further assume all quark masses to be equal to each other
\begin{equation}
	m_{f, 0} = m_0\,, \qquad f=1,\ldots,n_F\,.
\end{equation}
Therefore, the different quarks are indistinguishable and the
mixing between two different quark flavors cannot depend
on the flavors.

Defining the analogous operators of \eqn{eq:ops} for flowed fields, we
write
\begin{equation}
  \begin{split}
    \tcalo_{1,\mu\nu}(t,x) &\equiv
    G_{\mu\rho}^a(t,x)G_{\nu\rho}^a(t,x)\,,\\[10pt]
    \tcalo_{2,\mu\nu}(t,x) &\equiv
    \delta_{\mu\nu}G_{\rho\sigma}^a(t,x)G_{\rho\sigma}^a(t,x)\,,\\[10pt]
    \tcalo_{3_f,\mu\nu}(t,x) &\equiv Z_{\chi_f}
    \bar\chi_{f}(t,x)\left(\gamma_\mu\overleftrightarrow{\mathcal{D}}\!_\nu
    +\gamma_\nu\overleftrightarrow{\mathcal{D}}\!_\mu\right)\chi_{f}(t,x)\,,\\
    \tcalo_{4_f,\mu\nu}(t,x)
    &\equiv Z_{\chi_f} \delta_{\mu\nu}\bar\chi_{f}(t,x)
    \overleftrightarrow{\slashed{\mathcal{D}}}\chi_{f}(t,x)\,,
    \\
    \tcalo_{i,\mu\nu} &= \sum_{f=1}^{n_F}\tcalo_{i_f,\mu\nu}\,,\qquad i\in\{3,4\}\,,
    \label{eq:tops}
  \end{split}
\end{equation}
where $Z_{\chi_f}\equiv Z_\chi$ is the renormalization constant for the
flowed quark fields, and
\begin{equation}
  \begin{split}
    \overleftrightarrow{\mathcal{D}}\!_\mu = \partial_\mu -
    \overleftarrow{\partial}\!_\mu + 2B_\mu\,.
    \label{eq:leftrightcalD}
  \end{split}
\end{equation}
Since we have eliminated $\calo_{5,\mu\nu}$ from the set of operators by
using \eqn{eq:eom}, we do not need to include a flowed version of this
operator in \eqn{eq:tops}. Similar to the composite operators of regular
\qcd, we assume that the vacuum expectation values of the flowed
composite operators have been subtracted,
i.e.\ $\langle\tcalo_{i,\mu\nu}(t,x)\rangle\equiv 0\ \forall i$.

We can now use the expansion in small flow time\,\cite{Luscher:2011bx}
\begin{equation}
  \begin{split}
    \tcalo_{i,\mu\nu}(t,x) = \zeta_{ij}(t)\calo_{j,\mu\nu}(x) + \ldots\,,
    \label{eq:zetas}
  \end{split}
\end{equation}
to get a relation between flowed and regular \qcd\ operators. In
\eqn{eq:zetas}, and similarly in what follows, a sum $\sum_{j=1}^4$ is
understood. The ellipsis denotes terms that vanish as $t\to 0$ which
will be neglected throughout this paper. As discussed above, matrix
elements of the l.h.s.\ of this equation are finite after
renormalization of the \qcd\ parameters, while those of the regular
\qcd\ operators on the r.h.s.\ are in general divergent.  The mixing
matrix $\zeta_{ij}(t)$ will therefore be divergent as well.

Inverting \eqn{eq:zetas} and using it to re-express the regular
\qcd\ operators in the energy-momentum tensor in terms of flowed fields,
one arrives at
\begin{equation}
  \begin{split}
    T_{\mu\nu}(x) =  c_i (t) \tcalo_{i, \mu \nu} (t, x) \,,
    \label{eq:emtflow}
  \end{split}
\end{equation}
where
\begin{equation}
  \begin{split}
    c_i(t) \equiv \frac{1}{g_{0}^2} \left( \zeta_{1i}^{-1} (t)
		 - \frac{1}{4} \zeta_{2i}^{-1} (t) \right)
		 + \frac{1}{4} \zeta_{3i}^{-1} (t)
                         \,,\qquad i=1,\ldots,4\,.
    \label{eq:EMTGF}
	\end{split}
\end{equation}
Since matrix elements of the $\tcalo_i$ as well as the energy-momentum
tensor itself are finite (after mass and charge renormalization), the
universal coefficients $c_i(t)$ of \eqn{eq:EMTGF} are finite as
well.
In \citere{Makino:2014taa}, they have been calculated in perturbation
theory through \nlo\ \qcd. The goal of the current paper is to evaluate them
through \nnlo\ \qcd.
